\documentclass[11pt]{amsart}
\usepackage{amsmath,amssymb,amsthm,mathtools}
\usepackage{thmtools}
\usepackage{cleveref}

\newtheorem{theorem}{Theorem}[section]
\newtheorem{proposition}[theorem]{Proposition}
\newtheorem{lemma}[theorem]{Lemma}
\newtheorem{corollary}[theorem]{Corollary}

\theoremstyle{definition}
\newtheorem*{definition}{Definition}
\newtheorem*{question}{Question}
\newtheorem*{answer}{Answer}

\newcommand{\Tdk}{\mathcal{T}_{d,k-1}}


\begin{document}
\section{Preliminaries}\label{sec:prelim}
We state the elementary preliminaries required for the argument.

\subsection{Lucas theorem and its variants}
We record the elementary digit facts used throughout.

\begin{definition}[Carry-free addition]\label{def:carryfree}
For nonnegative integers $x_1,\ldots,x_r$, we write
\[
  x_1\oplus x_2\oplus\cdots\oplus x_r
\]
when the addition is carry-free in base-$p$.  Equivalently, if we have
\[
  x_i=\sum_{e\geq 0}x_{i,e}p^e
  \qquad {\text {\rm and}}\qquad 0\leq x_{i,e}\leq p-1,
\]
then we have
\[
  \sum_{i=1}^r x_{i,e}\leq p-1
  \qquad\text{for every }e.
\]
\end{definition}

The first tool is Lucas' theorem in its binomial form,
which provides a bridge between nonvanishing modulo $p$ and carry-free decompositions.

\begin{lemma}[Lucas' theorem]\label{lem:Lucas-binomial}
Let $m,n\geq 0$, and write
\[
  m=\sum m_ep^e\qquad {\text {\rm and}}\qquad n=\sum n_ep^e,
  \qquad 0\leq m_e,n_e\leq p-1.
\]
Then
\[
  \binom{n}{m}\equiv\prod_e\binom{n_e}{m_e}\pmod p.
\]
In particular, $\binom{n}{m}\not\equiv 0\pmod p$ if and only if $m_e\leq n_e$ for every $e$;
that is, if $m \oplus (n-m)$ is carry-free in base $p$.
\end{lemma}

We will mostly use the following multinomial form.

\begin{lemma}[Lucas' theorem, multinomial form]\label{lem:mult-Lucas}
Let $R=r_1+\cdots+r_m$.  Then
\[
  \binom{R}{r_1,\ldots,r_m}\not\equiv 0\pmod p
\]
if and only if $r_1 \oplus r_2 \oplus \dots \oplus r_m$ is carry-free in base-$p$.
\end{lemma}

\begin{proof}
Write the multinomial coefficient as a product of binomial coefficients:
\[
  \binom{R}{r_1,\ldots,r_m}
  =\binom{r_1+\cdots+r_m}{r_m}
   \binom{r_1+\cdots+r_{m-1}}{r_{m-1}}\cdots
   \binom{r_1+r_2}{r_2}.
\]
Each factor is nonzero modulo $p$ exactly when the indicated addition has no carry.
This is equivalent to saying that, at each base-$p$ digit position,
the digits of $r_1,\ldots,r_m$ add to at most $p-1$.
\end{proof}

\subsection{Finite field power sums}
The other elementary input is the standard finite-field power sum.
It is responsible for the divisibility conditions by $q-1$ that appear
throughout the later minimization problems.

\begin{lemma}[Finite-field power sums]\label{lem:field-sums}
For $m\geq 0$,
\[
  \sum_{x\in\mathbb{F}_q}x^m=
  \begin{cases}
    -1, & m>0\text{ and }(q-1)\mid m,\\
    0, & \text{otherwise}.
  \end{cases}
\]
Here we consider $0^0=1$ for $m=0$.
\end{lemma}

\begin{proof}
If $m=0$, the sum is $q$, which is $0$ in characteristic $p$.
If $m>0$, the zero term contributes $0$,
and the nonzero elements of $\mathbb{F}_q$ form a cyclic group of order $q-1$.
The sum of the $m$th powers in this cyclic group is $0$ unless $(q-1)\mid m$;
in that case it is instead $q-1=-1$ in $\mathbb{F}_q$.
\end{proof}


\section{Main proof}
Write a monic polynomial of degree $d$ as
\[
  a=t^d+\theta_1t^{d-1}+\cdots+\theta_d\qquad \text{and} \qquad \theta_i\in\mathbb{F}_q.
\]
Then
\[
  a^{-k}=t^{-dk}(1+\theta_1t^{-1}+\cdots+\theta_dt^{-d})^{-k}.
\]
Expanding the last factor gives
\begin{align*}
S_d(k)
&=t^{-dk}
  \sum_{\theta_1,\ldots,\theta_d\in\mathbb{F}_q}
  \sum_{m_1,\ldots,m_d\geq 0}
  \binom{-k}{m_1+\cdots+m_d}
  \binom{m_1+\cdots+m_d}{m_1,\ldots,m_d}\\
&\qquad\qquad\qquad\cdot
  \theta_1^{m_1}\cdots\theta_d^{m_d}
  t^{-(m_1+2m_2+\cdots+dm_d)}.
\end{align*}

By \Cref{lem:field-sums}, summing over $\theta_i\in\mathbb{F}_q$ kills every term except those with
\[
  m_i>0\qquad {\text {\rm and}}\qquad (q-1)\mid m_i.
\]
Next we turn attention to the binomial and multinomial coefficient.
Put $y \coloneq m_1 + \dots + m_d$.
The identity
\[
  \binom{-k}{y}=(-1)^y\binom{k+y-1}{y}
\]
shows, by \Cref{lem:Lucas-binomial},
that $\binom{-k}{y}$ is nonzero modulo $p$ if and only if $(k-1)+y$ is carry-free.
By \Cref{lem:mult-Lucas},
the multinomial coefficient is nonzero modulo $p$ if and only if $m_1+\cdots+m_d=y$ is carry-free.
These two carry-free conditions hold simultaneously if and only if the combined addition
\[
  (k-1)\oplus m_1\oplus\cdots\oplus m_d
\]
is carry-free, i.e., if $(m_1, \dots, m_d) \in \Tdk$.
\end{document}
